\documentclass[twoside]{article}
\usepackage[preprint]{aistats2027}
\usepackage[round,authoryear]{natbib}
\usepackage{amsmath,amssymb,amsthm}
\usepackage{mathtools}
\usepackage{microtype}
\usepackage{booktabs}
\usepackage{url}
\renewcommand{\bibname}{References}
\renewcommand{\bibsection}{\subsubsection*{\bibname}}
\bibliographystyle{apalike}
\newtheorem{theorem}{Theorem}
\newtheorem{lemma}[theorem]{Lemma}
\newtheorem{proposition}[theorem]{Proposition}
\newtheorem{corollary}[theorem]{Corollary}
\newtheorem{remark}[theorem]{Remark}
\newtheorem{definition}[theorem]{Definition}
\newcommand{\E}{\mathbb E}
\newcommand{\Pp}{\mathbb P}
\newcommand{\Q}{\mathcal Q}
\newcommand{\A}{\mathcal A}
\newcommand{\R}{\mathbb R}
\newcommand{\h}{\mathsf h}
\newcommand{\TV}{\mathsf{TV}}
\newcommand{\Pdim}{\operatorname{Pdim}}
\begin{document}
\runningauthor{}
\twocolumn[
\aistatstitle{Supplement: Near-Parametric Learning of Autoregressive Trajectories}
\aistatsauthor{Pahan Dewasurendra}
\aistatsaddress{Johns Hopkins University}
]

\appendix
\section{The Conditional $\rho$ Specialization}
\label{app:rho}

We give the reduction from the general results of \citet{baraud2018} to the unknown-design conditional $\rho$ theorem in the main paper. This also verifies that the unknown prompt marginal is absent from the algorithm.

Let $\mathcal X$ be a measurable prompt space, $\mathcal Z$ a finite trajectory space, and $\lambda$ counting measure on $\mathcal Z$. Let $\mu$ be the unknown prompt marginal. A conditional density $q$ is a measurable nonnegative function on $\mathcal X\times\mathcal Z$ satisfying $\sum_zq(z\mid x)=1$ for every $x$. It induces a joint probability
\[
 Q_q(dx,dz)=\mu(dx)q(z\mid x).
\]
The true joint probability has the disintegration $P(dx,dz)=\mu(dx)P_x(dz)$.

\begin{lemma}[Conditional Hellinger identity]
\label{lem:conditional-h}
For every conditional density $q$,
\[
 \h^2(P,Q_q)=\int_{\mathcal X}\h^2(P_x,q(\cdot\mid x))\,d\mu(x).
\]
\end{lemma}

\begin{proof}
Both probabilities are dominated by $\mu\otimes\lambda$. Expanding the definition and applying Tonelli gives
\begin{align*}
 \h^2(P,Q_q)
 &=\frac12\int\sum_z
 \left(\sqrt{P_x(z)}-\sqrt{q(z\mid x)}\right)^2d\mu(x)\\
 &=\int\h^2(P_x,q(\cdot\mid x))d\mu(x).
\end{align*}
\end{proof}

For $n$ independent observations, \citet{baraud2018} use the Hellinger-type distance that sums the $n$ marginal squared distances. Thus their product-space loss is $n\h^2(P,Q_q)$. Their Proposition 7 bounds the $\rho$-dimension of a VC-subgraph density class by
\begin{equation}
 D_n\leq C_0(V\wedge n)\left[1+\log_+(n/V)\right].
 \label{eq:rho-dimension}
\end{equation}
Their high-probability single-model inequality, stated in Equation (20) and normalized for i.i.d. density estimation in their Equation (29), then gives
\begin{equation}
 \h^2(P,Q_{\widehat q})
 \leq C_1\left[
 \inf_{q\in\Q}\h^2(P,Q_q)+\frac{D_n+\xi}{n}
 \right]
 \label{eq:source-rho}
\end{equation}
with probability at least $1-e^{-\xi}$. Their Section 8 proves that the same estimator applies to conditional densities without knowing $\mu$.

To see this directly, every candidate joint density relative to $\mu\otimes\lambda$ is $q(z\mid x)$. For two candidates,
\[
 \frac{dQ_{q'}}{dQ_q}(x,z)=\frac{q'(z\mid x)}{q(z\mid x)}.
\]
The prompt density cancels. Taking $\psi(r)=(r-1)/(r+1)$, the empirical comparison statistic is exactly
\[
 T_n(q,q')=\sum_{i=1}^n\psi\!\left(
 \sqrt{\frac{q'(Z_i\mid X_i)}{q(Z_i\mid X_i)}}
 \right).
\]
It can be evaluated from the sample and the candidate conditional densities alone.

\begin{proof}[Proof of the unknown-design conditional $\rho$ theorem]
Apply \eqref{eq:source-rho} and \eqref{eq:rho-dimension}, take $\xi=\log(1/\delta)$, and use Lemma~\ref{lem:conditional-h}. Proposition 8 of \citet{baraud2018} preserves the VC-subgraph index when passing from one-observation densities to their i.i.d. product representation. Their Theorem 3 permits a universally separable uncountable density class by reducing it to a fixed countable pointwise-dense subclass. This proves the displayed oracle inequality.

We finally verify the last proper-selection statement. Enumerate the countable subclass as $(q_j)_{j\geq1}$. Every $T_n(q_j,q_k)$ is measurable, and so are the countable supremum over $k$ and infimum over $j$. The approximate minimizer set in Definition 3 of \citet{baraud2018} is nonempty. Selecting its first element in the enumeration is measurable and is itself a $\rho$-estimator. It lies in the candidate subclass rather than only its Hellinger closure.
\end{proof}

\begin{lemma}[Solving the fixed-point rate]
\label{lem:fixed-point}
Let $V\geq1$, $\epsilon\in(0,1)$, and
\[
 n\geq\frac{C}{\epsilon}\left[V\log(C/\epsilon)+\log(1/\delta)\right]
\]
for a sufficiently large universal $C$. Then
\[
 \frac{(V\wedge n)[1+\log_+(n/V)]+\log(1/\delta)}{n}
 \leq c\epsilon
\]
for any prescribed universal $c>0$ after increasing $C$.
\end{lemma}

\begin{proof}
For sufficiently large $C$, the assumed inequality implies $n>V$. Put $r=n/V>1$. Dropping the nonnegative confidence term from the sample-size lower bound gives $r\geq C\epsilon^{-1}\log(C/\epsilon)$. Since $(1+\log r)/r$ is decreasing for $r\geq1$, substitution at this lower endpoint bounds it by $c\epsilon$ after increasing $C$. The confidence term is controlled directly by the second summand in the sample-size lower bound.
\end{proof}

\section{Algebraic Trajectory Complexity}
\label{app:algebra}

We record two variants of the trajectory-closure calculation.

\begin{lemma}[Pairwise ratio subgraphs]
\label{lem:ratio-vc}
Under the complexity-$(p,\Delta)$ assumption, if every transition probability is strictly positive, the class
\[
 \left\{(x,z,r):\frac{q_\theta(z\mid x)}{q_{\theta'}(z\mid x)}>r\right\},
 \qquad \theta,\theta'\in\Theta,
\]
has VC dimension less than
\[
 4p\log_2(16eM\Delta).
\]
\end{lemma}

\begin{proof}
Write $q_\theta=N_\theta/D_\theta$, where both trajectory polynomials have degree at most $M\Delta$. The ratio test is
\[
 N_\theta D_{\theta'}-rD_\theta N_{\theta'}>0.
\]
This is one polynomial inequality of degree at most $2M\Delta$ in the $2p$ parameters $(\theta,\theta')$. Applying Warren's sign-pattern bound exactly as in the proof of the trajectory-closure theorem gives the stated bound.
\end{proof}

This lemma explains directly why pairwise trajectory comparisons remain simple. The density-subgraph bound in the main paper is sharper and is the exact hypothesis used by Proposition 7 of \citet{baraud2018}.

\begin{lemma}[Observed bounded stopping]
\label{lem:stopping}
Suppose generation stops according to any rule whose decision at time $t$ is a deterministic function of the observed prefix and whose stopping time is at most $M$. If the stopped path includes its length, the stopped-trajectory density class obeys the same pseudo-dimension bound as the trajectory-closure theorem.
\end{lemma}

\begin{proof}
For a fixed stopped observation $(m,z_{1:m})$, the event that the stopping rule accepts this length is determined by $(x,z_{1:m})$ and introduces no parameter inequality. Its density is the product of the first $m\leq M$ transition probabilities, possibly multiplied by a fixed indicator. The proof of the trajectory-closure theorem applies with degree at most $m\Delta$.
\end{proof}

\subsection{Binary logistic delay line}

For a fixed path, let $Y_t=\operatorname{tail}_d(xz_{<t})$ and set $u^y=\prod_ju_j^{y_j}$. The complete density has the explicit form
\begin{equation}
 q_u(z\mid x)=
 \frac{\prod_{t:z_t=1}u^{Y_t}}
 {\prod_{t=1}^M(1+u^{Y_t})}.
 \label{eq:explicit-logistic-path}
\end{equation}
The numerator degree is at most $Md$, and the denominator degree is at most $Md$. No common denominator across unobserved paths is used.

The positive rational vectors $u\in\mathbb Q_{>0}^d$ are pointwise dense in $(0,\infty)^d$. For every fixed path, \eqref{eq:explicit-logistic-path} is continuous in $u$. Thus the universal-separability condition and the proper-selection step in the unknown-design conditional $\rho$ theorem both hold.

\subsection{Multiclass softmax}

Let token zero be the reference class. For $a\in\{1,\ldots,K-1\}$ and binary feature $y\in\{0,1\}^d$, define $u_a^y=\prod_{j=1}^du_{a,j}^{y_j}$. The conditional probabilities are
\begin{align*}
 g_u(a\mid y)&=\frac{u_a^y}{1+\sum_{b=1}^{K-1}u_b^y},\\
 g_u(0\mid y)&=\frac1{1+\sum_{b=1}^{K-1}u_b^y}.
\end{align*}
Each numerator and denominator has degree at most $d$ in $p=(K-1)d$ positive parameters. The trajectory-closure theorem therefore gives
\[
 V<2(K-1)d\log_2(8eMd).
\]
The bound changes only through the free parameter count when the state feature is any fixed binary map.

\section{Loss Conversions}
\label{app:loss}

\begin{lemma}[Coarsening and final-token loss]
\label{lem:coarsening}
Let $P,Q$ be probabilities on $\A^M$ and let $\phi$ be any measurable statistic. Then
\[
 \h^2(P\circ\phi^{-1},Q\circ\phi^{-1})\leq\h^2(P,Q).
\]
For binary final-token probabilities $p$ and $q$,
\[
 (p-q)^2\leq2\h^2(\operatorname{Bern}(p),\operatorname{Bern}(q)).
\]
Also $\TV(P,Q)\leq\sqrt2\h(P,Q)$.
\end{lemma}

\begin{proof}
The first claim is the data-processing inequality for Hellinger distance. It also follows from Cauchy--Schwarz after grouping the summands over each fiber of $\phi$. The remaining claims are the standard inequalities $\TV^2\leq2\h^2$ and $\TV\leq\sqrt2\h$ under the normalization used here.
\end{proof}

Applying Lemma~\ref{lem:coarsening} pointwise in $X$ proves the final-token claim in the near-parametric stochastic logistic CoT theorem. For trajectory sampling,
\begin{align*}
 \E_X\TV(P_X,Q_X)
 &\leq\sqrt2\E_X\h(P_X,Q_X)\\
 &\leq\sqrt{2\E_X\h^2(P_X,Q_X)}.
\end{align*}
Thus integrated squared Hellinger error $\epsilon^2/2$ implies expected trajectory total variation at most $\epsilon$.

\section{Minimax Lower Bound}
\label{app:lower}

We prove the minimax parameter lower-bound theorem in full. The argument applies to every output predictor, whether or not it corresponds to a logistic generator.

\begin{proof}[Proof of the minimax parameter lower bound]
First assume $d\geq8$. Let the prompt space contain $x_1,\ldots,x_d$ with $\operatorname{tail}_d(x_j)=e_j$, and let $X$ be uniform on these prompts. By the Varshamov--Gilbert bound, there is $\Omega\subseteq\{-1,+1\}^d$ with
\[
 \log|\Omega|\geq c_1d,
 \qquad
 d_H(\omega,\omega')\geq d/4
\]
for distinct $\omega,\omega'\in\Omega$.

Fix $\alpha\in(0,1/4]$ and let
\[
 \gamma=\log\frac{1/2+\alpha}{1/2-\alpha}.
\]
For $\omega\in\Omega$, set $w^\omega_j=\gamma\omega_j$. At horizon one,
\[
 f_\omega(x_j)=\sigma(w^\omega_j)=\frac12+\alpha\omega_j.
\]
Consequently,
\begin{equation}
 \|f_\omega-f_{\omega'}\|_{L_2(\mu)}^2
 =\frac{4\alpha^2}{d}d_H(\omega,\omega')
 \geq\alpha^2.
 \label{eq:packing-separation}
\end{equation}

Let $P_\omega$ be the law of one pair $(X,Z_1)$ under this target. For Bernoulli parameters in $[1/4,3/4]$,
\[
 \operatorname{KL}(\operatorname{Bern}(p)\|\operatorname{Bern}(q))
 \leq\frac{(p-q)^2}{q(1-q)}\leq22\alpha^2
\]
when $p,q\in\{1/2-\alpha,1/2+\alpha\}$. Averaging over the prompt gives
\begin{equation}
 \operatorname{KL}(P_\omega\|P_{\omega'})\leq22\alpha^2.
 \label{eq:packing-kl}
\end{equation}

Take $\alpha^2=16\epsilon$, with $\epsilon\leq1/256$. If an estimator $\widehat f$ has squared $L_2(\mu)$ error at most $\epsilon$, then its nearest element of $\{f_\omega:\omega\in\Omega\}$ is the true $f_\omega$. Indeed, the distance to the truth is at most $\sqrt\epsilon$, while \eqref{eq:packing-separation} gives pairwise distance at least $4\sqrt\epsilon$. Fano's inequality and \eqref{eq:packing-kl} show that this decoding succeeds with probability at most $2/3$ whenever $n\leq c_2d/\epsilon$. Hence no learner can guarantee error at most $\epsilon$ with success probability $3/4$ in this range.

For $d<8$, the same order up to a universal constant follows from the two hypotheses at one coordinate below. It remains to obtain the confidence dependence. Use one prompt and two targets with probabilities $1/2\pm\alpha$, now taking $\alpha=2\sqrt\epsilon$. An estimate with squared error at most $\epsilon$ identifies the sign. The sum of the two testing errors is at least
\[
 \begin{aligned}
 &\frac12\exp\left[-n\operatorname{KL}\left(
 \operatorname{Bern}(1/2+\alpha)\|\operatorname{Bern}(1/2-\alpha)
 \right)\right]\\
 &\hspace{28mm}\geq\frac12e^{-C n\epsilon}.
 \end{aligned}
\]
by the Bretagnolle--Huber inequality. Making both errors at most $\delta$ requires $n\geq c_3\log(1/\delta)/\epsilon$. Taking the larger of the dimension and confidence lower bounds, then changing the universal constant, proves the stated lower bound.
\end{proof}

\section{Sparse Model Selection}
\label{app:sparse}

For $S\subseteq[d]$, let $\Q_S$ be the logistic trajectory family with $w_j=0$ for $j\notin S$. Put $s=|S|$ and
\[
 V_s=\left\lceil2(s\vee1)\log_2(8eM(s\vee1))\right\rceil.
\]
The harmless use of $s\vee1$ includes the all-zero model.

Choose model weights
\begin{equation}
 \Delta(S)=\log\binom d{|S|}+2\log(|S|+1)+\log(\pi^2/6).
 \label{eq:support-penalty}
\end{equation}
They satisfy
\[
 \sum_{S\subseteq[d]}e^{-\Delta(S)}
 =\frac6{\pi^2}\sum_{s=0}^d\frac1{(s+1)^2}\leq1.
\]

\begin{theorem}[Sparse trajectory oracle]
\label{thm:sparse-full}
There is one penalized conditional $\rho$-estimator such that, with probability at least $1-\delta$,
\begin{align*}
 \E_X\h^2(P_X,\widehat q_X)
 \leq C\inf_{S\subseteq[d]}\bigg\{&\inf_{q\in\Q_S}\E_X\h^2(P_X,q_X)\\
 &+\frac{V_s[1+\log_+(n/V_s)]}{n}\\
 &+\frac{\Delta(S)+\log(1/\delta)}{n}\bigg\}.
\end{align*}
\end{theorem}

\begin{proof}
Each fixed-support class is universally separable and has pseudo-dimension at most $V_s$ by the trajectory-closure theorem. Apply the penalized model-selection theorem of \citet[Theorem 2]{baraud2018} with the Kraft weights \eqref{eq:support-penalty}. Its model complexity is bounded by Proposition 7 as in \eqref{eq:rho-dimension}. Divide the product-space Hellinger loss by $n$ and apply Lemma~\ref{lem:conditional-h}.
\end{proof}

If the truth has support size $s$, use $\Delta(S)\leq s\log(ed/s)+2\log(s+1)+C$. Lemma~\ref{lem:fixed-point} then yields the sparse rate displayed in the main paper up to secondary logarithms.

\section{Misspecification and Contamination}
\label{app:robustness}

The approximation term in the unknown-design conditional $\rho$ theorem already proves the misspecified claim. We add details for Huber contamination because the contaminating distribution may alter the prompt marginal.

Let $P$ be the clean joint law of a prompt and trajectory, and let
\[
 R=(1-\eta)P+\eta H
\]
be the observed law. Denote their prompt marginals by $\mu_P$ and $\mu_R$. Assume the clean conditional density is $q^\star\in\Q$. The candidate joint law used by the conditional estimator is
\[
 Q_R^\star(dx,dz)=\mu_R(dx)q^\star(z\mid x).
\]
Since $\TV(P,R)\leq\eta$ and marginalization contracts total variation,
\begin{align*}
 \h^2(P,R)&\leq\eta,\qquad
 \h^2(\mu_P,\mu_R)\leq\eta,\\
 \h^2(P,Q_R^\star)
 &=\h^2(\mu_Pq^\star,\mu_Rq^\star)
 =\h^2(\mu_P,\mu_R)\leq\eta.
\end{align*}
The squared triangle inequality gives $\h^2(R,Q_R^\star)\leq4\eta$. The unknown-design conditional $\rho$ theorem applied to the observed law therefore gives
\[
 \h^2(R,\mu_R\widehat q)\leq C(\eta+\mathfrak C_n)
\]
with high probability, where $\mathfrak C_n$ is the estimation term. Finally,
\begin{align*}
 \h(P,\mu_P\widehat q)
 \leq{}&\h(P,R)+\h(R,\mu_R\widehat q)\\
 &+\h(\mu_R\widehat q,\mu_P\widehat q),
\end{align*}
and the last term equals $\h(\mu_R,\mu_P)$. Squaring proves
\[
 \E_{X\sim\mu_P}\h^2(P_X,\widehat q_X)
 \leq C'(\eta+\mathfrak C_n).
\]
This is the contamination statement in the main paper.

\section{Why the e2e Calculation Differs}
\label{app:e2e}

For one observed CoT path, Equation~\eqref{eq:explicit-logistic-path} has one denominator product of degree at most $Md$. Under e2e supervision, the observed density is instead
\[
 q^{\mathrm{e2e}}_u(1\mid x)
 =\sum_{z_{1:M}:z_M=1}q_u(z_{1:M}\mid x).
\]
The summands have path-dependent denominators. The proof of \citet[Lemma 12]{doron2026stochastic} clears them using
\[
 D(u)=\prod_{y\in\{0,1\}^d}(1+u^y)^M,
\]
whose degree is $Md2^{d-1}$. The Goldberg--Jerrum bound then becomes
\[
 O\!\left(d\log(Md2^d)\right)=O(d^2+d\log M).
\]
Our observed-path closure theorem does not change this sum-before-threshold operation. Any improvement for e2e supervision requires additional structure beyond the present algebraic argument.

\bibliography{references}
\end{document}
